\odiseachapter{eolo}{La bolsa de los vientos}\label{cap:eolo}

Después de escapar del Cíclope, los aqueos llegan a Eolia, isla donde habita Eolo, el señor de los vientos, o, más literalmente, \gr{ταμίης ἀνέμων} (tamíēs anémōn). La palabra \gr{ταμίης} se traduce más bien como administrador o despensero. De hecho, el Eolo de la \emph{Odisea} no es un dios sino un mortal con funciones de mayordomo. Eso sí, goza de una posición acomodada en la isla de Eolia, donde lo pasa estupendamente junto a sus hijas e hijos, que no solo están juntos sino incestuosamente revueltos:

\begin{versocita}
Hasta que al fin la isla Eolia tocamos; en ella vivía\\
Eolo Hipótada, amado del cielo. Era aquella una isla\\
flotante y toda a redor por un muro de bronce ceñida\\
inquebrantable, y alzadas de roca que erguíanse lisas.\\
En su palacio doce hijos nacieron, y allí los tenía;\\
seis hijas tuvo, y seis hijos, y, entrados en flor de la vida,\\
a cada hijo le vino a dar él por esposa una hija.\\
Junto a su padre querido y su honrada madre, los días\\
ellos en fiesta los pasan, que siempre hay mil viandas servidas\\
y las paredes rezuman contento y buen husmo a cocina;\\
eso de día, que luego la noche a cada uno lo arrima\\
al lecho junto a su santa con los cobertores por cima.
\end{versocita}

Eolo ofrece hospitalidad y los navegantes aprovechan para tomarse un respiro:

\begin{versocita}
A su ciudad, pues, llegamos, y a aquellas estancias magníficas.\\
Un mes entero me tuvo él allí, que de todo quería\\
detalle: Ilión, el regreso de aqueos, las naves argivas…\\
Yo de eso todo, y a ley, por menudo el relato le hacía.
\end{versocita}

Cuando Odiseo le pide permiso para continuar su viaje, el simpático Eolo no solo se lo da, sino que le ofrece un regalo, como ya sabemos que exige la \emph{xenía}, cumplida aquí al pie de la letra. Eolo le regala a Ulises nada menos que la bolsa de los vientos. Es un tesoro que le permitirá llegar a su patria en un santiamén, pero, curiosamente, su magia exige no dejar que nadie se acerque a ella y, desde luego, no abrirla ---estamos, claro, ante otra versión del mito de Pandora---. Lo que ha hecho el señor de los vientos es aprisionar a todos sus díscolos súbditos en un saco muy bien cerrado, para que no molesten durante la travesía.

\begin{versocita}
Cuando le hablé de seguir ya camino y rogué despedida,\\
él no me dijo que no, y hasta me preparó la partida.\\
Tras desollar un buey nueveañero y hacerme una odrina,\\
tráfago en ella de vientos hinchados me ató que mugían,\\
pues despensero de vientos lo había hecho a él el Cronida\\
y en él estaba excitar o calmar al ciclón que quería.\\
Luego en la nave con cuerda de plata otros nudos le hacía\\
porque ni por las costuras soplara pizca de brisa.\\
Nos levantó un sostenido poniente al salir de la isla\\
para llevar a navíos y hombres, mas no llegaría\\
eso a cumplirse: nuestro desatino forjó nuestra ruina.
\end{versocita}

¿Qué insensatez van a cometer Ulises y los suyos esta vez? Una de juzgado de guardia:

\begin{versocita}
Hicimos rumbo muy bien, noche y día, por nueve jornadas\\
cuando en la décima se apareció como luz nuestra patria,\\
y hasta veíamos los que atendían allá sus fogatas.\\
Y de repente ahí se entró el dulce sueño en mis fuerzas cansadas,\\
que es que fui siempre a la escota del barco, y no quise dejarla\\
a hombre ninguno, por antes mi tierra poder alcanzarla.
\end{versocita}

Ponte un segundo en la piel de Odiseo y sus hombres, empática lectora, compasivo lector. Desde la cubierta del barco, al anochecer, estos hombres cansados, que se han pasado diez años peleando una guerra sucia y llevan ya largas semanas de viaje, pueden distinguir las hogueras que arden en su isla, pueden casi imaginarse a los pastores que quizás las hayan encendido para la trasnochada, se ven ya, literalmente, pisando la arena de sus playas. Es el momento en el que parece que, por fin, sus penas están a punto de terminar.

¿Por qué retrasa Odiseo la llegada? ¿Tan cansado está? O quizás, avispado como es, quiere llegar de día, con toda la gloria de un rey. No olvidemos que, en este momento de la travesía, todavía está al mando de numerosas naves, una flotilla con la que ---podemos imaginar--- planea llegar a plena luz del sol, anunciando su retorno, cargado de tesoros, a amigos y potenciales enemigos, que sin duda tomarán nota de sus numerosas huestes.

O quizás se dice a sí mismo que esa es su última noche de aventuras antes de regresar al hogar, a la vida doméstica, a la administración de sus tierras y sus ganados. Quizás le cuesta decidirse a desprenderse de su manto de héroe, a dar por concluidas sus aventuras y resignarse a ser solo otro hombre, \gr{Οὖτις}, Nadie.

O, simplemente, los dioses, siempre dispuestos a hacer sufrir a los mortales para divertirse, le aturden con un sopor al que no puede sobreponerse. Y Ulises se echa a dormir a las mismas puertas de su patria.

\begin{versocita}
Mis compañeros, entonces, cambiaban entre ellos palabras,\\
y se decían que yo me llevaba oro y plata a mi casa\\
que del magnánimo Hipótada Eolo, seguro, eran dádiva.\\
Y así uno entre ellos decía cruzando con otro miradas:
\stanzabreak
«Ay, hay que ver cómo quieren y cómo honran todos a este,\\
sea cual sea la tierra y ciudad a la que se acerque.\\
Muchos tesoros de Troya se trae en botín, y excelentes;\\
y ya lo ves, hecho el mismo viaje, a nosotros, su gente,\\
con nuestras manos vacías a casa llegando nos tiene.\\
Y encima ahora se lleva en amor y compaña el presente\\
de Eolo. Pronto, vayamos a ver lo que sea que encierre,\\
la tanta plata y el oro que el cuero, seguro, guarece».
\end{versocita}

A las mismas puertas de la patria, la traición. Nolan captura bien en su film la naturaleza levantisca de los compañeros del héroe, pero Homero es mucho más despiadado: los pinta ingratos, envidiosos e imprudentes. Y por supuesto, se buscan su propia ruina:

\begin{versocita}
Así decían. Venció el mal consejo de mis compañeros,\\
y desataron el odre, y aupáronse todos los vientos.\\
Un huracán que sobre ellos cayó los llevó mar adentro\\
llorando, porque la patria se les marchaba. Yo, en esto,\\
me desperté y en mi pecho sin tacha me vi discurriendo\\
si de la nave saltar y acabar de una vez en el piélago,\\
o soportarlo en silencio y seguir con los vivos más tiempo.\\
Y soporté, y me quedé, y me tendí, mi cabeza cubriendo\\
sobre las tablas, mientras arrastraba el ciclón de mal viento\\
a la isla Eolia las naves de nuevo, y mis hombres gimiendo.
\end{versocita}

Los versos son estremecedores. El engañador engañado, el estafador estafado por su propia gente. Y de paso, otro mito deslizándose entre los versos, el de Sísifo, que empuja su piedra montaña arriba solo para verla rodar por la ladera cuando ya casi había alcanzado la cumbre, igual que Ulises ve alejarse Ítaca cuando ya se imaginaba fondeando en su puerto. La desesperación es tan grande que considera lanzarse al mar, quitarse la vida. ¿Qué le dolió más, afligida lectora, compungido lector? ¿La patria que se alejaba o la traición de sus amigos?

De vuelta a Eolia, Odiseo suplica clemencia, humillándose ante Eolo:

\begin{versocita}
Desembarcamos allí para hacer de agua dulce el abasto,\\
y por las lévedas naves comiendo enseguida ya estábamos.\\
En cuanto todos de yanta y bebida la gana amansamos,\\
a un compañero conmigo tomé y a otro más como heraldo\\
y fui al glorioso palacio de Eolo; a él allí festejando\\
me lo encontré con su esposa y su prole. No entré en el palacio,\\
que, sin cruzar el umbral, ante el mismo portón nos sentamos.\\
Y se asombraron aquellos al vernos y nos preguntaron:
\stanzabreak
«¿Cómo, Odiseo, tú aquí? ¿Qué demonio se te ha echado encima?\\
Te despedimos de gana y de fe porque así llegarías\\
a esa tu patria y tu casa, o adonde quisieras venida».
\stanzabreak
Eso dijeron, y yo les hablé desde mi desconsuelo:\\
«Daño me dio la mi mala compaña, y daño el mal sueño;\\
mas remediádmelo, amigos, que está en vosotros hacerlo».
\end{versocita}

Pero esta vez no hay suerte. Eolo considera que su antiguo huésped está maldito y lo echa a patadas:

\begin{versocita}
«Sal de la isla ya mismo, que vivo no lo hay más maléfico;\\
yo permitido no tengo amparar ni mandar de regreso\\
a hombre que sea de dioses benditos aborrecimiento.\\
Sal, que si has vuelto tú aquí, como odiado del cielo lo has hecho».
\end{versocita}

Y los aqueos, con el corazón «sombrirroto», salen a remo, ya que no sopla una brizna de viento:

\begin{versocita}
Así diciendo, me echó de su casa. Y yo en hondo sollozo.\\
De allí otra vez navegamos con el corazón sombrirroto,\\
mientras roía a los hombres su espíritu el remo penoso\\
por necedad, que ya allí no asomaba la escolta de un soplo.
\end{versocita}

Por necedad, todo por necedad. A veces da la impresión de que la estupidez humana es el tema central de la \emph{Odisea}.

Navegan seis días y al séptimo llegan a la isla de los lestrigones. Lo que Odiseo cuenta entonces a los feacios anuncia ya un oscuro presagio.

\begin{versocita}
Cuando a su espléndido puerto llegamos, que por escarpados\\
roquedos y de continuo ceñido va a lado y a lado, y\\
dos puntas bravas, enfrente una de otra, se ven estrechando\\
la embocadura, dejando al final solo un paso delgado,\\
todos pasaron por él sus barcos belquilibrados\\
y los dejaron ahí dentro en el cuéncavo puerto amarrados\\
flanco con flanco, que grande o pequeña se estaba a resguardo\\
aquel lugar de la ola, en un luminoso mar calmo.\\
Solo yo el negro navío afuera dejé sin entrarlo, y\\
cable le eché en un extremo del puerto contra unos peñascos.
\end{versocita}

¿Por qué Odiseo no entra en el puerto como las demás naves\footnote{Todavía cuenta con los doce navíos que partieron de Troya, ya que no ha perdido ninguna nave ni en la incursión con los cícones ni en el encuentro con Polifemo.}? ¿Teme algo o es su naturaleza desconfiada y siempre alerta? La ensenada es a la vez un puerto seguro y una trampa, una de tantas trampas de las que escapa el héroe, mitad por astucia, mitad por suerte.

Odiseo decide explorar y las cosas se ponen feas enseguida. En una escena que Nolan reproduce con gran economía narrativa, los aqueos encuentran a una niña ---bien robusta, nos informa el canto--- que resulta ser la hija de Antífates, el lestrigón. Esta los lleva a casa de su padre, que lo primero que hace es merendarse a uno de los viajeros, mientras los otros huyen despavoridos.

\begin{versocita}
Se hizo el almuerzo con uno al que, visto y no visto, echó mano;\\
los otros dos, escapándose al vuelo, a las naves llegaron.\\
Mas levantó aquel un grito que por la ciudad fue escuchado\\
y congregó a lestrigones robustos de todos los lados,\\
innumerables, y más a gigantes semejos que a humanos.
\end{versocita}

Es decir, los aqueos están viviendo la segunda entrega de la serie «Cíclopes», con la diferencia de que estos gigantes, al contrario que los paisanos de Polifemo, están organizados y son muy numerosos. Esta vez no hay astucia que valga y sigue una carnicería espantosa:

\begin{versocita}
Como quien tira un guijarro, lanzaban pedruscos de lo alto, y\\
casi a la vez se enalzaba en las naves estrépito malo\\
de hombres muriendo y de barcos con ellos allí destrozados.\\
Y como a peces los ensartaban para un yantar agrio.
\end{versocita}

Observa, atenta lectora, despierto lector, la mezcla de ironía y brutal justicia que dispensa el Poeta. Estos aqueos que han saqueado Troya, atacado sin motivo a los cícones y cegado al cíclope, son ahora diezmados por hombres todavía más brutales que ellos, unos gigantones ---magníficamente reimaginados por Nolan--- que los ensartan como peces y les hunden las naves sin provocación alguna. Quien a hierro mata…, parece decirnos el bardo. Solo la nave de Odiseo escapa a la carnicería.

\begin{versocita}
Mientras adentro del puerto honduroso hacían ellos estrago,\\
tirando yo de la espada buída que va a mi costado,\\
cable enseguida con ella corté del navío proiblavo.\\
Y presuroso a mis hombres moví y ordené sin descanso\\
jalar de remos para sortear ese tanto quebranto;\\
y sacudían el mar todos ellos, de muerte espantados.\\
A la maralta escapó felizmente de aquellos peñascos\\
mi nave; todas las otras allí su final encontraron.
\end{versocita}

Pero el héroe ya no es el capitán de una flota, sino un superviviente que dispone de un único barco. Si estimamos unos cincuenta tripulantes por nave, acaba de perder más de quinientos hombres en la refriega y solo le quedan unas pocas decenas. El depredador ahora es presa, las tornas han cambiado y Homero cierra el episodio con las mismas líneas que cerraba el canto dedicado al cíclope:

\begin{versocita}
De allí otra vez navegamos, el pecho a dolores transido,\\
salvos de muerte, mas faltos de los compañeros queridos.
\end{versocita}
